\section{Preliminaries}\label{sec:preliminaries}

The proof uses two Fourier expansions. Characters of affine rows encode
candidate labels, whereas characters of Boolean table entries detect the
influences to which Majority Is Stablest applies. We introduce both
conventions, then state the exact affine-game input.

All vector spaces over $\F$ are finite dimensional. All expectations on a
finite vector space or Boolean cube use the uniform probability measure,
unless another distribution is specified. All Fourier coefficients use
this probability normalization. An $L^2$ norm always integrates over the
variables displayed in its subscript; variables declared fixed are not
integrated.

\subsection{Affine rows and labels}

Let $\mathcal A$ be a nonempty affine space over $\F$, with a fixed affine
coordinate system, and set
\[
 H=\Aff(\mathcal A,\F).
\]
This is a vector space, containing the constant form $\one$.
For $q\in H^*$ define $\chi_q(h)=(-1)^{q(h)}$.
Evaluation identifies $\mathcal A$ with
\begin{equation}\label{eq:labels-dual}
 \{q\in H^*:q(\one)=1\}.
\end{equation}
Indeed, if $h(a)=h_0+\sum_i h_i a_i$, a functional taking value one on
$\one$ has the unique form $q(h)=h_0+\sum_i h_i a_i$.
We call precisely these functionals \emph{labels}; a general nonzero
functional in $H^*$ need not be a label.

For $f:H^k\to\mathbb R$ and $C=(q_1,\ldots,q_k)\in(H^*)^k$, put
\[
 \widehat f(C)=\E_R f(R)\prod_{i=1}^k\chi_{q_i}(R_i),\qquad
 |C|=\#\{i:q_i\ne0\}.
\]
Parseval gives $\sum_C\widehat f(C)^2=\E_R f(R)^2$.
The same notation is used for functions with additional variables, which
remain fixed during the indicated Fourier expansion.

For $0\le r\le1$, let $T_r$ be the row-noise operator: independently in each
coordinate, retain the row with probability $r$ and otherwise replace it
by an independent uniform row. Then
\begin{equation}\label{eq:row-noise}
 \widehat{T_rf}(C)=r^{|C|}\widehat f(C).
\end{equation}
This operator preserves the range $[-1,1]$. It is also convolution by an
independent increment $E\in H^k$, whose coordinates are zero with
probability $r$ and independent uniform rows with probability $1-r$.
Convolution commutes with every fixed translation of the rows.

\subsection{Boolean Fourier analysis}

For a finite set $\Omega$ and $h:\bits^\Omega\to\mathbb R$, write
\[
 h(x)=\sum_{S\subseteq\Omega}\widehat h(S)\prod_{a\in S}x_a.
\]
Its influence and degree-$K$ influence at $a\in\Omega$ are
\[
 \Inf_a(h)=\sum_{S\ni a}\widehat h(S)^2,\qquad
 \Inf_a^{\le K}(h)=\sum_{\substack{S\ni a\\|S|\le K}}
 \widehat h(S)^2.
\]
Let $P_{\le K}$ be the orthogonal projection onto Fourier degree at most
$K$. In particular,
\begin{equation}\label{eq:degree-influence}
 \sum_{a\in\Omega}\Inf_a(P_{\le K}h)
 \le K\norm{h}_2^2.
\end{equation}
For $t$-correlated uniform bits $x,y$, independently coupled in each
coordinate,
\begin{equation}\label{eq:stability}
 \Stab_t(h)=\E[h(x)h(y)]
 =\sum_{S\subseteq\Omega}t^{|S|}\widehat h(S)^2.
\end{equation}
Thus stability is nonnegative for $t\ge0$. If $h\in[-1,1]$, it is at most
one.

\begin{lemma}[Majority Is Stablest with a degree cutoff]\label{lem:mis}
Fix $0<t<1$ and $\xi>0$. There are an integer $K\ge1$ and $\tau>0$,
depending only on $t,\xi$, such that every
$h:\bits^\Omega\to[-1,1]$ with $\E h=0$ satisfies
\[
 \bigl(\Inf_a^{\le K}(h)<\tau\text{ for all }a\in\Omega\bigr)
 \quad\Longrightarrow\quad
 \Stab_t(h)\le B(t)+\xi.
\]
\end{lemma}

\begin{proof}
The untruncated assertion follows from Majority Is Stablest
\cite[Theorem~4.4]{MOO10}. To match its $[0,1]$ convention, put
$g=(h+1)/2$: then $\E g=1/2$, each nonconstant Fourier coefficient is
halved, and $\Stab_t(h)=4\Stab_t(g)-1$. The Gaussian halfspace value
for $g$ is $1/4+\arcsin(t)/(2\pi)$, which becomes $B(t)$ in our
normalization; see \cite[Theorem~4.1 and Corollary~4.3]{MOO10} and
\cite{Borell85}. Here is the degree-cutoff reduction to the untruncated
assertion. For a cube-noise
parameter $u<1$, put $h_u=T_u h$, using ordinary bit noise in this
paragraph. Range and mean are preserved. Uniformly over $h$ of norm at
most one,
\[
 0\le\Stab_t(h)-\Stab_t(h_u)
 \le\sup_{j\ge0}t^j(1-u^{2j})\longrightarrow0
 \quad(u\uparrow1).
\]
The convergence follows by first choosing a finite degree cutoff to bound
the tail $t^j$, and then taking $u$ close to one on the remaining degrees.
Choose $u$ so that this error is at most $\xi/2$. Let $\delta>0$ be the
untruncated influence threshold for error $\xi/2$. Choose $K$ with
$u^{2(K+1)}\le\delta/2$, and take $\tau=\delta/2$. Then
\[
 \Inf_a(h_u)
 \le\Inf_a^{\le K}(h)+u^{2(K+1)}\norm{h}_2^2
 <\delta.
\]
Apply the untruncated theorem to $h_u$ and add the stability error.
\end{proof}

\subsection{Affine projection games}

A weighted bipartite projection game consists of finite sets $X,Y$, a
rational probability distribution $\lambda$ on a finite edge set
$E$ with endpoint maps to $X$ and $Y$ (parallel edges are allowed), finite alphabets,
and a projection $\pi_e$ for each edge. Its value is
\[
 \Val(\mathcal G)=
 \max_{\ell_X,\ell_Y}\Pr_{e=(x,y)\sim\lambda}
 [\pi_e(\ell_X(x))=\ell_Y(y)].
\]
In this paper the alphabets are $\F^{n+1}$ and $\F^n$, and every
$\pi_e:\F^{n+1}\to\F^n$ is surjective affine.
\Cref{prop:affine-games} supplies NP-hard gaps with this structure.
We discard endpoints with zero marginal probability whenever conditional
edge distributions are used.

The starting theorem has the quantifier order
\[
 \forall b\in(0,1)\quad \exists n=n(b)\ge1\quad
 \forall a\in(0,1-b).
\]
The precise promise is given below; its derivation from
the published construction is in \Cref{sec:affine-games}.

\begin{proposition}[Affine 2-to-1 Games Theorem]
\label{prop:affine-games}
For every constant $b\in(0,1)$ there is an integer $n\geq 1$ such that,
for every constant $a\in(0,1-b)$, the following promise problem is NP-hard.
Its input is a
weighted bipartite game with vertex sets $X,Y$, respective label sets
$\F^{n+1},\F^n$, and a surjective affine map
\[
  \pi_e:\F^{n+1}\longrightarrow\F^n
\]
on each edge $e=(x,y)$.  An edge is satisfied when
$\pi_e(\lambda_X(x))=\lambda_Y(y)$.  The promise cases are
\[
  \mathrm{YES}:\quad \Val(\mathcal G)\geq 1-a,
  \qquad
  \mathrm{NO}:\quad \Val(\mathcal G)\leq b.
\]
The integer $n$ depends on $b$ alone.  For fixed $a,b$, the reduction is
deterministic and has polynomial running time and output size.
\end{proposition}

A weighted graph is likewise specified by a rational distribution on
unordered edges, allowing loops. A cut assigns signs to its vertices and
satisfies an edge when its signs differ. Its value is its satisfaction
probability. Loops have zero cut weight. The conversion in
\Cref{lem:unweighted} keeps a common scale even when loops are deleted.
